% =============================================================================
% Project : THE PHILOSOPHY OF VIABILITY
% Book    : Book I — Ontology of Viability
% Book ID : VS.B01
% Part    : Part I — Identity and Manifestation
% Part ID : VS.B01.P01
% Chapter : Chapter 08 — Comparative Philosophical Examination
% Chapter ID: VS.B01.P01.C08
%
% File :
%   2_MAINMATTER/PART_I_IDENTITY_AND_MANIFESTATION/
%   CHAPTER_08_COMPARATIVE_PHILOSOPHICAL_EXAMINATION/
%   sections/section_12_al_farabi.tex
%
% Section :
%   Section 12 — Al-Farabi
%
% Structural Identifier :
%   VS.B01.P01.C08.S12
%
% Purpose :
%   To reconstruct Al-Farabi's account of metaphysics, Being,
%   Unity, quiddity, the First Existent, causal dependence,
%   emanation, the cosmological hierarchy, and the relation
%   between cosmic order and the virtuous city; and to examine
%   the significance and limits of these positions in relation
%   to the ontology of Being and Identity developed in Part I.
%
% Method :
%   Historical and textual reconstruction;
%   distinction between primary claims and later interpretation;
%   analysis of disputed textual and doctrinal relations;
%   symmetrical comparative examination;
%   strongest challenge and considered response;
%   bilateral assessment of explanatory costs;
%   substitution test and explicit claim-status classification.
%
% Foundational Restrictions :
%   1. Being and Identity remain analytically distinguishable
%      and ontologically inseparable.
%   2. Analytical Distinction, Ontological Inseparability,
%      and Absence of Ontological Priority remain separate
%      foundational commitments.
%   3. Identity is not equated with quiddity, determination,
%      unity, singularity, individuation, numerical identity,
%      invariance, structure, composition, or manifestation.
%   4. The First Existent is not identified with Being as such,
%      First Determination, or any proposed realm of Book I.
%   5. Emanation is not equated with Expression or Manifestation.
%   6. The Active Intellect is not identified with Light.
%   7. The hierarchy of intellects and celestial spheres is
%      not identified with the proposed domains and hubs
%      of the Ontology of Viability.
%   8. Al-Farabi's account of intellect and human cognition
%      is not made a foundation of Book I's ontology.
%   9. The political hierarchy of the virtuous city is examined
%      historically and comparatively, not endorsed as a
%      normative commitment of this book.
%
% Claim Classification :
%   Functional or logical form and epistemic or justificatory
%   status are classified independently.
%
% Bibliographical Basis :
%   AlFarabiVirtuousCity
%   MennAlFarabiMetaphysicsSEP2026
%   GermannAlFarabiSocietyReligionSEP2021
%
% Asset Inventory :
%   TL01 — Comparative claims and ontological limits.
%   DG01 — Schematic account of the emanative cosmology.
%   FG   — Not produced; no independent figure is required.
%   IL   — Deferred; a separate illustration would risk
%          conflating the political city with the cosmic order.
%
% Technical Note :
%   Paragraph anchors are provided through standard hyperref
%   labels. No new internal-reference macro is assumed.
%
% =============================================================================


\section{Al-Farabi}
\label{sec:ch08_s12}

\index[name]{Al-Farabi}
\index[concept]{First Existent}
\index[concept]{First Cause}
\index[concept]{Being}
\index[concept]{Existence}
\index[concept]{Unity}
\index[concept]{Quiddity}
\index[concept]{Emanation}
\index[concept]{Causal Dependence}
\index[concept]{Separate Intellects}
\index[concept]{Active Intellect}
\index[concept]{Celestial Spheres}
\index[concept]{Virtuous City}
\index[concept]{Metaphysics}
\index[concept]{Determination}
\index[concept]{Identity}
\index[concept]{Manifestation}


% =============================================================================
% 1. SCOPE AND COMPARATIVE ORIENTATION
% =============================================================================

\subsection{Scope and Comparative Orientation}
\label{subsec:ch08-farabi-scope}

% Structural ID: VS.B01.P01.C08.S12.SS01.P001
\phantomsection\label{para:VS.B01.P01.C08.S12.SS01.P001}

Al-Farabi is an important interlocutor for the ontology of
Being and Identity because his philosophical work addresses
several problems that cannot be reduced to the description
of isolated entities. His investigations concern the
meanings of Being and Unity, the relation between a thing
and its quiddity, the dependence of existents upon a first
principle, and the ordered plurality of beings. His
metaphysics is therefore relevant not only to the question
of the ultimate source of reality but also to the question
of how an existent is specified and distinguished from
what it is not.

% Structural ID: VS.B01.P01.C08.S12.SS01.P002
\phantomsection\label{para:VS.B01.P01.C08.S12.SS01.P002}

The examination must distinguish two related but different
areas of his work. The first concerns general metaphysical
questions: the senses of Being, existence, Unity, quiddity,
and the scientific investigation of their causes. These
questions are especially important in the
\textit{Book of Letters} and the treatise commonly known
as \textit{On the One and Unity}. The second concerns
the First Existent, its attributes, causal dependence,
emanation, and the cosmological order, particularly in
\textit{Al-Farabi on the Perfect State} and
\textit{The Political Regime}.
\citep[see especially secs.~3--4]{MennAlFarabiMetaphysicsSEP2026}

% Structural ID: VS.B01.P01.C08.S12.SS01.P003
\phantomsection\label{para:VS.B01.P01.C08.S12.SS01.P003}

The central comparative question is whether Al-Farabi's
account of the specification and unity of existents
provides an explanation of what makes an existent the
existent that it is, or whether it addresses a different
problem: the ways in which Being and Unity are predicated,
the way a thing is circumscribed by its quiddity, and
the dependence of a plurality of beings upon a first cause.
These questions are closely related, but their relation
must not be assumed to be one of conceptual identity.

% Structural ID: VS.B01.P01.C08.S12.SS01.P004
\phantomsection\label{para:VS.B01.P01.C08.S12.SS01.P004}

The principal thesis of this section is that Al-Farabi
offers a substantial analysis of Being, Unity, quiddity,
and causal order, but does not thereby provide a complete
account of Identity in the sense developed in Part I.
His account places pressure on the relation between
Identity and Determination and requires a more precise
examination of Unity. Yet neither the First Existent,
the unity attributed to it, nor the circumscription of
a thing by its quiddity can be substituted without
argument for Identity as a distinct ontological problem.

% Structural ID: VS.B01.P01.C08.S12.SS01.P005
\phantomsection\label{para:VS.B01.P01.C08.S12.SS01.P005}

This comparison does not seek to translate Al-Farabi's
conceptual vocabulary into the terminology of the
Philosophy of Viability. The First Existent is not
identified with Being as such; quiddity is not simply
renamed Identity; Unity is not treated as Singularity;
and emanation is not treated as another expression
for Manifestation. Each proposed relation requires
an examination of conceptual function, explanatory
scope, and inferential role.


% =============================================================================
% 2. HISTORICAL AND TEXTUAL LOCATION
% =============================================================================

\subsection{Historical and Textual Location}
\label{subsec:ch08-farabi-historical}

% Structural ID: VS.B01.P01.C08.S12.SS02.P001
\phantomsection\label{para:VS.B01.P01.C08.S12.SS02.P001}

Abu Nasr al-Farabi was a philosopher of the tenth
century whose work engaged extensively with Greek
philosophy and the intellectual traditions of the
Islamic world. His philosophical investigations
include logic, metaphysics, the classification of
the sciences, psychology, political philosophy,
and music. The breadth of this work is relevant
because his metaphysical claims are connected to
his accounts of human perfection, intellectual
activity, and political association. Nevertheless,
the presence of these subjects within a wider
philosophical project does not make them
interchangeable or establish that every work
presents an identical doctrinal system.

% Structural ID: VS.B01.P01.C08.S12.SS02.P002
\phantomsection\label{para:VS.B01.P01.C08.S12.SS02.P002}

The principal primary source used here is
\textit{Al-Farabi on the Perfect State}, Richard
Walzer's translation of the work commonly known
as \textit{Mabadi' Ara' Ahl al-Madina al-Fadila}.
The translated title emphasizes the political
dimension of the work, but its contents extend
well beyond political organization. Its opening
metaphysical account concerns the First Existent,
its attributes, the origin of other existents,
and the order of the cosmos. Its later discussions
connect this account with the human being and
the virtuous city.
\citep[pp.~56--89, 229--233]{AlFarabiVirtuousCity}

% Structural ID: VS.B01.P01.C08.S12.SS02.P003
\phantomsection\label{para:VS.B01.P01.C08.S12.SS02.P003}

The \textit{Book of Letters} and
\textit{On the One and Unity} are particularly
important for the reconstruction of Al-Farabi's
general metaphysics. They address questions
that cannot be adequately reconstructed from
the cosmological passages of the \textit{Perfect
State} alone. The analysis of Being and Unity
in these works provides part of the conceptual
background for the arguments concerning the
First Existent and its attributes. The relation
between these treatises and the more explicitly
theological and cosmological works must,
however, be established carefully rather than
assumed to constitute a single uninterrupted
exposition.
\citep[secs.~3.1--3.6]{MennAlFarabiMetaphysicsSEP2026}

% Structural ID: VS.B01.P01.C08.S12.SS02.P004
\phantomsection\label{para:VS.B01.P01.C08.S12.SS02.P004}

The chronological order and compositional history
of Al-Farabi's works are not sufficiently settled
to justify treating every difference between
them as a straightforward stage in the development
of a single doctrine. Some tensions may reflect
different philosophical purposes, different
textual contexts, or developments in his thought.
The present comparison therefore distinguishes
claims made in a particular work from broader
reconstructions of Al-Farabi's philosophy.
\citep[secs.~1 and 4]{MennAlFarabiMetaphysicsSEP2026}

% Structural ID: VS.B01.P01.C08.S12.SS02.P005
\phantomsection\label{para:VS.B01.P01.C08.S12.SS02.P005}

A further historical restriction concerns the
relation between Al-Farabi and Avicenna. Both
philosophers examine Being, quiddity, causal
dependence, and the First Existent or Necessary
Principle. This continuity does not establish
that Al-Farabi already possesses the complete
technical distinction between essence and
existence developed in Avicenna's metaphysics.
The later formulation must not be projected
backward merely because related questions
appear in both systems. The transition to
Avicenna in the following section will examine
that difference directly.


% =============================================================================
% 3. METAPHYSICS AS A SCIENCE
% =============================================================================

\subsection{Metaphysics as a Science}
\label{subsec:ch08-farabi-metaphysics-science}

% Structural ID: VS.B01.P01.C08.S12.SS03.P001
\phantomsection\label{para:VS.B01.P01.C08.S12.SS03.P001}

Al-Farabi's metaphysical project is not confined
to the construction of a theological account
of God. It includes an attempt to reconstruct
metaphysics as a science concerned with the
most universal concepts and principles of
reality. This project is connected to his
interpretation of Aristotle and to his effort
to clarify the proper subject, method, and
questions of metaphysical inquiry. Being,
Unity, essence, cause, and the First Existent
belong to this wider investigation.
\citep[secs.~1--2]{MennAlFarabiMetaphysicsSEP2026}

% Structural ID: VS.B01.P01.C08.S12.SS03.P002
\phantomsection\label{para:VS.B01.P01.C08.S12.SS03.P002}

This methodological concern is significant
for the present comparison. Al-Farabi does
not assume that a philosophical term has
one unambiguous meaning simply because it
is used repeatedly. He investigates how
terms acquire different meanings in ordinary
language, in the sciences, and in philosophical
argument. He is especially attentive to the
possibility that grammatical form and
translation may conceal distinctions in
logical or philosophical function.
\citep[secs.~3.1--3.3]{MennAlFarabiMetaphysicsSEP2026}

% Structural ID: VS.B01.P01.C08.S12.SS03.P003
\phantomsection\label{para:VS.B01.P01.C08.S12.SS03.P003}

The attention to language must not be mistaken
for the claim that Being is constituted by
language. Al-Farabi's analysis of terminology
is intended to clarify the concepts and
questions through which metaphysics investigates
reality. The fact that a term has several
meanings does not establish that the realities
to which those meanings refer are produced
by linguistic convention. His methodological
analysis and his ontological commitments
must therefore be kept distinct.

% Structural ID: VS.B01.P01.C08.S12.SS03.P004
\phantomsection\label{para:VS.B01.P01.C08.S12.SS03.P004}

This distinction is relevant to the ontology
of Book I. An analysis of the meanings of
Being may expose conceptual confusions,
but it does not by itself settle the
ontological relation between Being and
Identity. Conversely, the ontological
independence of what exists from description
does not make conceptual analysis dispensable.
The task is to distinguish the reality
being investigated from the terms and
arguments through which that investigation
is conducted.


% =============================================================================
% 4. THE SENSES OF BEING AND EXISTENCE
% =============================================================================

\subsection{The Senses of Being and Existence}
\label{subsec:ch08-farabi-being-existence}

% Structural ID: VS.B01.P01.C08.S12.SS04.P001
\phantomsection\label{para:VS.B01.P01.C08.S12.SS04.P001}

Al-Farabi's analysis of Being distinguishes
several senses of the term commonly rendered
as \textit{mawjud}, as well as corresponding
uses of \textit{wujud}. These distinctions
are important because they prevent the
word ``being'' from being treated as though
it always performs one and the same
philosophical function. His account includes
Being as divided among the categories,
Being in the sense of truth, and Being
in the sense of what is circumscribed by
a quiddity outside the soul.
\citep[sec.~3.4]{MennAlFarabiMetaphysicsSEP2026}

% Structural ID: VS.B01.P01.C08.S12.SS04.P002
\phantomsection\label{para:VS.B01.P01.C08.S12.SS04.P002}

Being as divided among the categories
concerns the ways in which something is
said to be according to the categories
of Aristotelian metaphysics. This use
does not exhaust the meaning of Being.
It identifies one way in which an existent
is classified within a metaphysical
account of substance and the different
kinds of predication. The category-based
use must therefore be distinguished
from the more universal questions
concerning Being and its causes.

% Structural ID: VS.B01.P01.C08.S12.SS04.P003
\phantomsection\label{para:VS.B01.P01.C08.S12.SS04.P003}

Being in the sense of truth concerns
the relation between what is understood
or asserted and what obtains outside
the soul. In this sense, the truth of
a judgment or concept depends upon
whether what is understood corresponds
to something outside the soul. The
corresponding use of existence is
therefore relational: it concerns the
truth of what is thought or asserted.
This is not the same as claiming that
the existence of an external entity
depends upon its being thought.
\citep[sec.~3.4]{MennAlFarabiMetaphysicsSEP2026}

% Structural ID: VS.B01.P01.C08.S12.SS04.P004
\phantomsection\label{para:VS.B01.P01.C08.S12.SS04.P004}

Being in the sense of being circumscribed
by a quiddity outside the soul concerns
the existent as having a quiddity that
specifies what it is. This sense is not
reducible to the truth of a proposition
about the existent. Something may have
a quiddity outside the soul whether or
not it has been represented in the soul.
The distinction therefore prevents the
epistemic or truth-related use of Being
from being confused with the ontological
status of the existent itself.
\citep[sec.~3.4]{MennAlFarabiMetaphysicsSEP2026}

% Structural ID: VS.B01.P01.C08.S12.SS04.P005
\phantomsection\label{para:VS.B01.P01.C08.S12.SS04.P005}

The distinctions among these senses
must not be simplified into a general
thesis that existence is either only
a mental concept or an additional
entity attached to every essence.
Al-Farabi's treatment of \textit{wujud}
is more differentiated. The relation
between a being and its existence
depends, in part, upon the kind of
quiddity under consideration and
upon the sense in which Being is
predicated. In particular, his account
does not simply reproduce the later
Avicennian formulation of the real
distinction between essence and
existence in possible beings.
\citep[sec.~3.4]{MennAlFarabiMetaphysicsSEP2026}

% Structural ID: VS.B01.P01.C08.S12.SS04.P006
\phantomsection\label{para:VS.B01.P01.C08.S12.SS04.P006}

The comparative importance of this
analysis lies in its insistence that
the question ``what is it?'' cannot
be collapsed into the question
``is it?'' or into the question
``is the statement about it true?''
These questions may be connected
within metaphysical inquiry, but
their conceptual functions differ.
The distinction is relevant to Book I
because it requires a careful account
of Being without reducing Being to
a single mode of predication or
to the conditions of its description.


% =============================================================================
% 5. UNITY AND THE MANY SENSES OF ONE
% =============================================================================

\subsection{Unity and the Many Senses of One}
\label{subsec:ch08-farabi-unity}

% Structural ID: VS.B01.P01.C08.S12.SS05.P001
\phantomsection\label{para:VS.B01.P01.C08.S12.SS05.P001}

Al-Farabi's analysis of Unity is especially
important because it rejects the assumption
that every use of ``one'' refers to the
same kind of unity. One may refer to
numerical unity, continuity, indivisibility,
the absence of a partner, or the
circumscription of something by an
attribute or quiddity. These senses
do not all presuppose the same relation
between one and many.
\citep[sec.~3.6]{MennAlFarabiMetaphysicsSEP2026}

% Structural ID: VS.B01.P01.C08.S12.SS05.P002
\phantomsection\label{para:VS.B01.P01.C08.S12.SS05.P002}

A thing that is one in number may be
opposed to a plurality of numerically
distinct things. A thing that is one
as continuous may be opposed to a
discrete plurality. A thing that is
one as indivisible may be opposed
to what is divided. A thing that is
one because it has no partner within
its species may be opposed to what
has another of the same species.
These meanings cannot be reduced
to one general opposition between
Unity and Plurality.

% Structural ID: VS.B01.P01.C08.S12.SS05.P003
\phantomsection\label{para:VS.B01.P01.C08.S12.SS05.P003}

Al-Farabi also distinguishes the
plurality that arises from a one
from the plurality that is opposed
to that one. Several points may
arise from the unity of a point,
but the plurality of points is not
the same as the continuous extension
that stands opposite to a point
considered as indivisible. Likewise,
several entities may each be unique
within their respective kinds without
their plurality being the contrary
of the uniqueness attributed to
each of them. The one and the many
therefore require analysis according
to the respect in which each term
is being used.
\citep[sec.~3.6]{MennAlFarabiMetaphysicsSEP2026}

% Structural ID: VS.B01.P01.C08.S12.SS05.P004
\phantomsection\label{para:VS.B01.P01.C08.S12.SS05.P004}

This analysis is significant for the
Ontology of Viability. Unity cannot
be treated as synonymous with
Singularity, Aggregate, Composition,
or the unity of a source. Nor can
the plurality arising from a unity
be automatically treated as the
plurality opposed to it. The sense
of Unity must be specified together
with its subject, respect, and
ontological function. Al-Farabi's
analysis strengthens the need for
these distinctions, even though
his own classification does not
supply the complete conceptual
architecture adopted in Book I.

% Structural ID: VS.B01.P01.C08.S12.SS05.P005
\phantomsection\label{para:VS.B01.P01.C08.S12.SS05.P005}

A particularly important sense of
Unity concerns a thing as circumscribed
by its quiddity. In this account,
the quiddity does not merely describe
a thing after it has already been
constituted as a separate existent.
It is involved in specifying the
thing as distinct from what it is not.
Al-Farabi's analysis therefore connects
a particular sense of Unity with
the determinate character of Being.
\citep[sec.~3.6]{MennAlFarabiMetaphysicsSEP2026}

% Structural ID: VS.B01.P01.C08.S12.SS05.P006
\phantomsection\label{para:VS.B01.P01.C08.S12.SS05.P006}

This connection must not be generalized
into the claim that every kind of
Unity is identical with quiddity,
or that Unity is identical with
Identity. Al-Farabi's argument concerns
specified senses of Unity and their
relation to Being. Its significance
for Book I is comparative: it reveals
that an adequate ontology must
distinguish the unity of an existent,
the specification of its quiddity,
and the Identity of the existent
without presuming that one term
can replace the others.


% =============================================================================
% 6. QUIDDITY AND CIRCUMSCRIPTION
% =============================================================================

\subsection{Quiddity and Circumscription}
\label{subsec:ch08-farabi-quiddity}

% Structural ID: VS.B01.P01.C08.S12.SS06.P001
\phantomsection\label{para:VS.B01.P01.C08.S12.SS06.P001}

The notion of circumscription is central
to Al-Farabi's account of one and
Being. A thing may be circumscribed
by a spatial boundary or place, but
circumscription may also concern
the way in which a thing is specified
by an attribute or by its quiddity.
The latter sense is not a spatial
boundary. It concerns the determination
of what the thing is and how it is
distinguished from what it is not.
\citep[sec.~3.6]{MennAlFarabiMetaphysicsSEP2026}

% Structural ID: VS.B01.P01.C08.S12.SS06.P002
\phantomsection\label{para:VS.B01.P01.C08.S12.SS06.P002}

Circumscription by quiddity is not
merely the assignment of a verbal
description. It concerns the way
in which the existent is specified
by what it is. The distinction between
one thing and another is therefore
not exhausted by their being assigned
different names or recognized under
different descriptions. The quiddity
has an ontological role within
Al-Farabi's account of the being
and unity of things.

% Structural ID: VS.B01.P01.C08.S12.SS06.P003
\phantomsection\label{para:VS.B01.P01.C08.S12.SS06.P003}

The analysis also raises the question
of individual specification. A shared
species or generic quiddity does not
necessarily exhaust the way in which
one individual is distinguished
from another. Al-Farabi's discussion
of more complete circumscription
allows the possibility that individual
beings are specified in ways not
captured by a definition common
to all members of a species. This
interpretation is important, but
its precise metaphysical implications
remain dependent upon the difficult
details of his account.
\citep[sec.~3.6]{MennAlFarabiMetaphysicsSEP2026}

% Structural ID: VS.B01.P01.C08.S12.SS06.P004
\phantomsection\label{para:VS.B01.P01.C08.S12.SS06.P004}

The relation between quiddity and
Identity must nevertheless remain
open to critical analysis. Quiddity
concerns what a thing is and the
way in which it is specified.
Identity concerns the existent as
the existent that it is. These
questions may intersect, but their
intersection does not establish
that Identity is reducible to
quiddity or that quiddity exhausts
the ontological constitution of
an existent.

% Structural ID: VS.B01.P01.C08.S12.SS06.P005
\phantomsection\label{para:VS.B01.P01.C08.S12.SS06.P005}

The same restriction applies to
Determination. Circumscription by
quiddity may be relevant to the
determinate character of an existent,
but it does not establish that
Identity is identical with Determination.
Nor does it establish that every
determination associated with an
existent is constitutive of its
Identity. The question of which
determinations participate in
constituting Identity remains a
distinct ontological problem within
Part I.


% =============================================================================
% 7. THE FIRST EXISTENT
% =============================================================================

\subsection{The First Existent}
\label{subsec:ch08-farabi-first-existent}

% Structural ID: VS.B01.P01.C08.S12.SS07.P001
\phantomsection\label{para:VS.B01.P01.C08.S12.SS07.P001}

In the metaphysical account of the
\textit{Perfect State}, Al-Farabi
begins from the First Existent.
The First is not introduced as one
member among other beings within
an already constituted collection.
It occupies a distinctive position
as the first cause of existence
for the other beings whose existence
does not depend upon human will
and choice. Its priority is therefore
causal and metaphysical, not merely
chronological.
\citep[pp.~56--89]{AlFarabiVirtuousCity}

% Structural ID: VS.B01.P01.C08.S12.SS07.P002
\phantomsection\label{para:VS.B01.P01.C08.S12.SS07.P002}

The First Existent is described as
uncaused, independent, simple,
and without the composition that
would make it dependent upon
constituent parts. It is not a body
and is not constituted through
matter and form. Its priority is
connected with the claim that
nothing more basic is required
to explain its existence. The
metaphysical attributes attributed
to the First are therefore intended
to express its independence from
the kinds of causes and compositions
that characterize dependent beings.
\citep[pp.~56--74]{AlFarabiVirtuousCity}

% Structural ID: VS.B01.P01.C08.S12.SS07.P003
\phantomsection\label{para:VS.B01.P01.C08.S12.SS07.P003}

The argument from the First's
independence to its simplicity
and uniqueness is not a merely
descriptive classification. It is
an attempt to establish that a
first principle cannot possess
the kinds of composition that
would make its existence dependent
upon more basic constituents.
The transition from simplicity
to uniqueness, however, involves
substantive inferential steps.
Those steps must be examined
rather than accepted merely
because the conclusion is part
of the system's architecture.
\citep[sec.~4]{MennAlFarabiMetaphysicsSEP2026}

% Structural ID: VS.B01.P01.C08.S12.SS07.P004
\phantomsection\label{para:VS.B01.P01.C08.S12.SS07.P004}

The First Existent is also described
through Unity. Its unity is not
treated as an accidental property
added to an otherwise constituted
subject. The First is one in a
manner connected with its own
quiddity and simplicity. This
account is distinctive because
the unity attributed to the First
is not simply numerical singularity
within a wider class of comparable
entities. It expresses the absence
of a partner and of internal
composition within the first
principle.
\citep[pp.~66--78]{AlFarabiVirtuousCity}

% Structural ID: VS.B01.P01.C08.S12.SS07.P005
\phantomsection\label{para:VS.B01.P01.C08.S12.SS07.P005}

The First Existent must not,
however, be identified with
Being as such or with Identity
as developed in Book I. Al-Farabi's
First is a determinate metaphysical
principle within an account of
causal dependence and the order
of existents. Being in its most
universal sense and Identity
as the ontological fact that an
existent is the existent it is
are not thereby exhausted by
the description of this first
principle. Its uniqueness cannot
be generalized into a theory
of the unity or Identity of
every existent.


% =============================================================================
% 8. CAUSAL DEPENDENCE AND EMANATION
% =============================================================================

\subsection{Causal Dependence and Emanation}
\label{subsec:ch08-farabi-emanation}

% Structural ID: VS.B01.P01.C08.S12.SS08.P001
\phantomsection\label{para:VS.B01.P01.C08.S12.SS08.P001}

Al-Farabi explains the dependence
of other existents upon the First
through an account of emanation.
The term commonly rendered as
emanation refers to the derivation
of the existence of other beings
from the First. The existence of
what proceeds from the First does
not add to the First's perfection,
nor does the First require an
instrument or an additional
condition within itself in order
to produce an effect.
\citep[pp.~88--92]{AlFarabiVirtuousCity}

% Structural ID: VS.B01.P01.C08.S12.SS08.P002
\phantomsection\label{para:VS.B01.P01.C08.S12.SS08.P002}

This account is intended to preserve
the independence and simplicity
of the First while explaining
the existence of a plurality of
dependent beings. The First is
not altered by the existence of
what proceeds from it. Its causal
activity does not require the
addition of a new property or
the acquisition of a perfection
that it previously lacked. The
relation is therefore asymmetrical:
the dependent beings require
the First, while the First does
not require them in the same
manner.
\citep[pp.~88--95]{AlFarabiVirtuousCity}

% Structural ID: VS.B01.P01.C08.S12.SS08.P003
\phantomsection\label{para:VS.B01.P01.C08.S12.SS08.P003}

Emanation must not be confused
with a temporal sequence in which
the First exists at one time and
subsequently begins to produce
other beings. The metaphysical
priority of the First concerns
causal dependence, not simply
an earlier position in time.
The account is also distinct
from the ordinary model of
production in which an agent
changes through acting upon
material or acquires a new
condition through the production
of an effect.

% Structural ID: VS.B01.P01.C08.S12.SS08.P004
\phantomsection\label{para:VS.B01.P01.C08.S12.SS08.P004}

A significant interpretive difficulty
concerns the relation between
the First's existence and the
existence of other beings. Al-Farabi
asserts that the other beings
proceed from the First, but the
argument that the First's existence
alone entails the existence of
other beings is not presented
with the same explicitness as
the arguments concerning the
First's simplicity and attributes.
The distinction between establishing
that the First is a cause and
explaining how it causes other
beings is therefore important.
\citep[sec.~4]{MennAlFarabiMetaphysicsSEP2026}

% Structural ID: VS.B01.P01.C08.S12.SS08.P005
\phantomsection\label{para:VS.B01.P01.C08.S12.SS08.P005}

For the Ontology of Viability,
the central comparative consequence
is that causal dependence does
not by itself provide a complete
account of the Identity of the
dependent existent. To establish
that an existent derives its
existence from a first cause
is not yet to explain what
makes that existent the existent
that it is, how it is distinguished
from other dependent existents,
or which determinations are
constitutive of its Identity.
The causal question and the
question of Identity must remain
analytically distinguishable.


% =============================================================================
% 9. THE COSMOLOGICAL ORDER
% =============================================================================

\subsection{The Cosmological Order}
\label{subsec:ch08-farabi-cosmology}

% Structural ID: VS.B01.P01.C08.S12.SS09.P001
\phantomsection\label{para:VS.B01.P01.C08.S12.SS09.P001}

Al-Farabi's account of emanation
extends into a hierarchical
cosmological order. The First
Existent is followed by separate
intellects associated with the
celestial order, while the
subsequent structure includes
celestial spheres and the
Active Intellect, which governs
the sublunary domain. This
hierarchy is not merely a
classification of objects.
It is an account of causal
and explanatory dependence
among different orders of
existent beings.
\citep[sec.~4]{MennAlFarabiMetaphysicsSEP2026}

% Structural ID: VS.B01.P01.C08.S12.SS09.P002
\phantomsection\label{para:VS.B01.P01.C08.S12.SS09.P002}

The familiar presentation of
ten intellects and nine celestial
spheres is a useful summary of
one influential reconstruction
of Al-Farabi's cosmology.
It must not, however, be
treated as an unqualified
description that makes the
arrangement identical in every
work. The \textit{Perfect State}
and the \textit{Political Regime}
do not present precisely the
same account of the immediate
effects proceeding from the
First. The number and arrangement
of the secondary intellects
must therefore be stated with
reference to the particular
text being examined.
\citep[sec.~4]{MennAlFarabiMetaphysicsSEP2026}

% Structural ID: VS.B01.P01.C08.S12.SS09.P003
\phantomsection\label{para:VS.B01.P01.C08.S12.SS09.P003}

The cosmological hierarchy is
significant because it connects
causal dependence with a
differentiated plurality of
beings. The plurality is not
treated as an unordered
collection, nor is every existent
assigned the same causal role.
Different levels have different
relations to the First and
different functions within
the cosmological account.
This structure is one of the
reasons why Al-Farabi's
metaphysics cannot be reduced
to a simple opposition between
one source and an undifferentiated
many.

% Structural ID: VS.B01.P01.C08.S12.SS09.P004
\phantomsection\label{para:VS.B01.P01.C08.S12.SS09.P004}

The diagram below represents
only the broad causal architecture
of this cosmological account.
It does not claim that every
treatise presents an identical
sequence, nor does it imply
that the stages are temporal
events. The Active Intellect
is shown separately to clarify
its role in the sublunary order;
it remains part of the series
of separate intellects in the
relevant cosmological scheme.

% =============================================================================
% Project : THE PHILOSOPHY OF VIABILITY
% Book    : Book I — Ontology of Viability
% Book ID : VS.B01
% Part    : Part I — Identity and Manifestation
% Part ID : VS.B01.P01
% Chapter : Chapter 08 — Comparative Philosophical Examination
% Chapter ID : VS.B01.P01.C08
%
% File :
%   4_ASSETS/diagrams/
%   VS.B01.P01.C08.S12.SS09.P004.DG01.al_farabi_emanative_cosmological_order.tex
%
% Section : Section 12 — Al-Farabi
%
% Asset :
%   Diagram DG01 — Al-Farabi's Emanative Cosmological Order
%
% Structural Identifier :
%   VS.B01.P01.C08.S12.SS09.P004.DG01
%
% Purpose :
%   To represent the broad causal architecture of Al-Farabi's
%   emanative cosmology and distinguish the First Existent,
%   the secondary intellects, the celestial order, the Active
%   Intellect, and the sublunary domain.
%
% Intellectual Benefit :
%   Clarifies the broad relations of causal dependence while
%   avoiding an unsupported one-to-one mapping of every
%   intellect to a celestial sphere across all relevant texts.
%
% Aesthetic / Experiential Benefit :
%   Provides a compact visual synthesis of a complex cosmological
%   structure.
%
% Representational Limits :
%   1. The diagram is schematic, not exhaustive.
%   2. It does not represent a temporal sequence.
%   3. The arrangement of immediate effects differs between
%      the Perfect State and the Political Regime.
%   4. The Active Intellect is shown separately to clarify
%      its function; it belongs to the relevant series of
%      separate intellects.
%   5. The diagram is not a representation of the domains
%      or hubs proposed in Book I.
%   6. The Active Intellect is not identified with Light.
%
% Dependencies :
%   TikZ; positioning library; arrows.meta. (no other libraries required;
%   all widths scale with \linewidth).
%
% =============================================================================

\begin{figure}[htbp]
\centering

\begin{tikzpicture}[
    font=\small,
    node distance=0.95cm,
    box/.style={
        rounded corners=4pt,
        align=center,
        inner sep=6pt,
        minimum height=1.2cm
    },
    first/.style={
        box,
        draw=orange!75!black,
        fill=yellow!15,
        line width=1.0pt,
        text width=0.50\linewidth
    },
    intellect/.style={
        box,
        draw=blue!75!black,
        fill=blue!8,
        line width=0.8pt,
        text width=0.50\linewidth
    },
    active/.style={
        box,
        draw=green!55!black,
        fill=green!9,
        line width=0.9pt,
        text width=0.50\linewidth
    },
    celestial/.style={
        box,
        draw=blue!65!black,
        fill=gray!8,
        line width=0.7pt,
        text width=0.26\linewidth
    },
    sublunary/.style={
        box,
        draw=blue!65!black,
        fill=gray!15,
        line width=0.8pt,
        text width=0.50\linewidth
    },
    note/.style={
        draw=gray!65,
        fill=gray!5,
        rounded corners=3pt,
        text width=0.80\linewidth,
        align=left,
        font=\footnotesize,
        inner sep=6pt
    },
    relation/.style={
        draw=blue!65!black,
        line width=0.85pt,
        -{Stealth[length=2mm]}
    },
    explanatory/.style={
        draw=gray!65,
        line width=0.7pt,
        dashed,
        -{Stealth[length=1.8mm]}
    },
    elabel/.style={
        font=\scriptsize,
        align=left,
        inner sep=2pt
    }
]

% ---------------------------------------------------------------------------
% Main vertical axis: First Existent -> Secondary Intellects -> Active Intellect
% ---------------------------------------------------------------------------

\node[first] (first) {
    \textbf{The First Existent / First Cause}\\
    Uncaused, independent, and simple
};

\node[intellect, below=of first] (secondary) {
    \textbf{Secondary Separate Intellects}\\
    An ordered series of intellectual principles
};

\node[active, below=of secondary] (active) {
    \textbf{Active Intellect}\\
    Terminal separate intellect in the familiar sublunary scheme
};

\node[sublunary, below=1.15cm of active] (sublunary) {
    \textbf{Sublunary Domain}\\
    Material beings, forms, souls, generation, and corruption
};

% ---------------------------------------------------------------------------
% Celestial Order (beside the series, outside it)
% ---------------------------------------------------------------------------

\node[celestial, right=1.7cm of secondary] (spheres) {
    \textbf{Celestial Order}\\
    Spheres and their associated order
};

% ---------------------------------------------------------------------------
% Frame: the Active Intellect belongs to the series of separate intellects
% ---------------------------------------------------------------------------

\draw[blue!60!black, densely dotted, line width=0.6pt, rounded corners=6pt]
    ([xshift=-3.5mm, yshift=9mm]secondary.north west) rectangle
    ([xshift= 3.5mm, yshift=-3.5mm]active.south east);

\node[font=\scriptsize\itshape, text=blue!60!black, anchor=north west,
      inner sep=2pt]
    at ([xshift=-3.5mm, yshift=9mm]secondary.north west)
    {Series of separate intellects};

% ---------------------------------------------------------------------------
% Relations
% ---------------------------------------------------------------------------

\draw[relation] (first.south) -- 
    node[elabel, right, xshift=2pt] {Causal dependence}
    (secondary.north);

\draw[relation] (secondary.south) -- 
    node[elabel, right, xshift=2pt] {Continued ordered dependence}
    (active.north);

\draw[explanatory] (secondary.east) -- 
    node[elabel, above, align=center] {associated\\with}
    (spheres.west);

\draw[relation] (active.south) -- 
    node[elabel, right, xshift=2pt] {Governing and formative relation}
    (sublunary.north);

% ---------------------------------------------------------------------------
% Interpretive limits
% ---------------------------------------------------------------------------

\node[note, anchor=north west]
    (limits) at ([xshift=-3.5mm, yshift=-7mm]sublunary.south west) {
    \textbf{Interpretive limits.} Schematic only: no uniform one-to-one
    correspondence between each separate intellect and each celestial
    sphere is claimed across all treatises. The Active Intellect is
    displayed apart for clarity, not as a principle outside the intellectual
    series. Relations are metaphysical and causal, not temporal stages.
};

\end{tikzpicture}

\caption{Schematic representation of the broad causal architecture of
Al-Farabi's emanative cosmology. The First Existent is the first cause of
the ordered series of existents. The secondary intellects are associated
with the celestial order, while the Active Intellect, as the terminal member
of the series, stands in a governing and formative relation to the sublunary
domain. The diagram does not imply that the relevant treatises present an
identical arrangement.}

\label{fig:VS.B01.P01.C08.S12.SS09.P004.DG01}

\end{figure}

% =============================================================================
% END OF DIAGRAM DG01
% =============================================================================

% Structural ID: VS.B01.P01.C08.S12.SS09.P005
\phantomsection\label{para:VS.B01.P01.C08.S12.SS09.P005}

This cosmological hierarchy
must not be identified with
the proposed domains and
intermediate hubs of Book I.
The two arrangements have
different philosophical purposes
and different explanatory
commitments. A hierarchy of
separate intellects and celestial
spheres is a feature of Al-Farabi's
cosmology; the proposed
architecture of Book I is not
a reproduction of that hierarchy.
Any comparison between them
must identify the relevant
structural resemblance without
transferring the causal or
theological commitments of
one system into the other.


% =============================================================================
% 10. THE ACTIVE INTELLECT
% =============================================================================

\subsection{The Active Intellect}
\label{subsec:ch08-farabi-active-intellect}

% Structural ID: VS.B01.P01.C08.S12.SS10.P001
\phantomsection\label{para:VS.B01.P01.C08.S12.SS10.P001}

The Active Intellect occupies
a distinctive position in
Al-Farabi's cosmological account.
It is not simply a human faculty
or an abstract name for the
capacity to reason. It belongs
to the order of separate
intellects and is associated
with the sublunary domain.
Its role connects the
cosmological hierarchy with
the processes through which
forms and intellectual activity
are explained.
\citep[sec.~4]{MennAlFarabiMetaphysicsSEP2026}

% Structural ID: VS.B01.P01.C08.S12.SS10.P002
\phantomsection\label{para:VS.B01.P01.C08.S12.SS10.P002}

The Active Intellect also has
an important role in Al-Farabi's
account of human cognition.
It is associated with the
actualization of intellectual
capacities and with the
acquisition of intelligible
forms. This role connects
his cosmology with his
psychology and epistemology.
Those connections are relevant
to a complete reconstruction
of his philosophical system,
but they do not make the
theory of human cognition
a foundation for the ontology
of Book I.
\citep[secs.~2.1--2.2]{GermannAlFarabiSocietyReligionSEP2021}

% Structural ID: VS.B01.P01.C08.S12.SS10.P003
\phantomsection\label{para:VS.B01.P01.C08.S12.SS10.P003}

The Active Intellect must not
be identified with Light in
the Philosophy of Viability.
The former is a separate
intellectual principle within
a particular cosmological
and psychological system.
Light remains an independent
ontological problem in Book I.
A similarity in explanatory
language or in the role of
mediation is insufficient
to establish conceptual
equivalence between them.

% Structural ID: VS.B01.P01.C08.S12.SS10.P004
\phantomsection\label{para:VS.B01.P01.C08.S12.SS10.P004}

The comparative relevance of
the Active Intellect is therefore
restricted. It demonstrates
how Al-Farabi connects a
hierarchical cosmology with
the explanation of sublunary
forms and intellectual activity.
It does not establish that
Book I requires a separate
intellectual substance, nor
does it provide a definition
of Manifestation, Expression,
or Light within the present
ontology.


% =============================================================================
% 11. THE VIRTUOUS CITY AND THE COSMIC ORDER
% =============================================================================

\subsection{The Virtuous City and the Cosmic Order}
\label{subsec:ch08-farabi-virtuous-city}

% Structural ID: VS.B01.P01.C08.S12.SS11.P001
\phantomsection\label{para:VS.B01.P01.C08.S12.SS11.P001}

Al-Farabi's political philosophy
is connected to his metaphysical
and cosmological account through
an analogy between the ordered
cosmos, the human organism,
and the virtuous city. The city
is examined as an association
whose members cooperate toward
a common end. Its organization
is compared with the differentiated
parts of an organism and with
the ordered structure of the
cosmos. This analogy gives
political association a place
within the wider philosophical
account of human perfection.
\citep[pp.~229--233]{AlFarabiVirtuousCity}

% Structural ID: VS.B01.P01.C08.S12.SS11.P002
\phantomsection\label{para:VS.B01.P01.C08.S12.SS11.P002}

The virtuous city is not defined
merely by the provision of
material necessities or by
the maintenance of political
order. Its purpose is connected
with the conditions through
which human beings may attain
their highest perfection and
happiness. Cooperation,
education, intellectual
development, and political
direction are therefore
related within Al-Farabi's
account of the excellent
society.
\citep[secs.~1--2]{GermannAlFarabiSocietyReligionSEP2021}

% Structural ID: VS.B01.P01.C08.S12.SS11.P003
\phantomsection\label{para:VS.B01.P01.C08.S12.SS11.P003}

The analogy between city and
cosmos is philosophically
important, but it must not
be interpreted as proof that
every political hierarchy
is justified by the existence
of a cosmic hierarchy. The
analogy expresses a proposed
correspondence between
differentiated structures and
their ordering principles.
Whether the political
consequences follow from
the metaphysical account
requires a further argument.
The descriptive reconstruction
of Al-Farabi's position does
not constitute an endorsement
of its political implications.

% Structural ID: VS.B01.P01.C08.S12.SS11.P004
\phantomsection\label{para:VS.B01.P01.C08.S12.SS11.P004}

For Book I, the principal
comparative lesson is that
an ordered whole may contain
differentiated parts without
reducing those parts to
interchangeable units. This
is relevant to the examination
of plurality, composition,
and relational organization.
Nevertheless, the political
function of a citizen, the
causal role of a celestial
body, and the Identity of
an existent are different
questions. A functional
position within a whole
does not by itself determine
the Identity of the entity
occupying that position.


% =============================================================================
% 12. INTERPRETIVE DISPUTES AND TEXTUAL LIMITS
% =============================================================================

\subsection{Interpretive Disputes and Textual Limits}
\label{subsec:ch08-farabi-interpretive-disputes}

% Structural ID: VS.B01.P01.C08.S12.SS12.P001
\phantomsection\label{para:VS.B01.P01.C08.S12.SS12.P001}

The reconstruction of Al-Farabi's
metaphysics faces several
interpretive difficulties.
His works do not all present
the same subject in the same
manner, and the relation
between his general metaphysics,
his philosophical theology,
his cosmology, and his political
philosophy requires analysis.
A unified account may be
possible, but it must be
established through textual
and argumentative examination
rather than assumed in advance.
\citep[secs.~1--4]{MennAlFarabiMetaphysicsSEP2026}

% Structural ID: VS.B01.P01.C08.S12.SS12.P002
\phantomsection\label{para:VS.B01.P01.C08.S12.SS12.P002}

One dispute concerns the relation
between the First Existent and
the subsequent intellectual
order. The \textit{Perfect State}
and the \textit{Political Regime}
do not offer identical descriptions
of the immediate effects proceeding
from the First. Consequently,
the familiar scheme of ten
intellects and nine spheres
must be treated as a reconstruction
whose details require attention
to the work under examination.
\citep[sec.~4]{MennAlFarabiMetaphysicsSEP2026}

% Structural ID: VS.B01.P01.C08.S12.SS12.P003
\phantomsection\label{para:VS.B01.P01.C08.S12.SS12.P003}

A second difficulty concerns
the strength of the arguments
through which Al-Farabi derives
the attributes of the First.
The transitions from causal
priority to simplicity, from
simplicity to uniqueness, and
from immateriality to the
First's status as intellect
involve inferential steps
whose adequacy has been
questioned. The presence
of a conclusion within the
system must therefore be
distinguished from the
demonstrative sufficiency
of the argument offered
for it.
\citep[sec.~4]{MennAlFarabiMetaphysicsSEP2026}

% Structural ID: VS.B01.P01.C08.S12.SS12.P004
\phantomsection\label{para:VS.B01.P01.C08.S12.SS12.P004}

A third difficulty concerns
the interpretation of the
cosmological account in
relation to political philosophy.
The analogy between cosmic
and political order may be
read as an explanatory
correspondence, a normative
model, or a combination
of the two. The exact force
of this relation should not
be reduced to a decorative
comparison, but neither
should every political
conclusion be treated as
automatically established
by the cosmological hierarchy.
\citep[sec.~3]{GermannAlFarabiSocietyReligionSEP2021}

% Structural ID: VS.B01.P01.C08.S12.SS12.P005
\phantomsection\label{para:VS.B01.P01.C08.S12.SS12.P005}

These disputes limit the degree
of certainty with which a
single comprehensive system
can be attributed to Al-Farabi.
The present section therefore
distinguishes claims supported
by the primary text from
interpretive reconstructions
and comparative conclusions.
Where the available evidence
does not establish a unique
reading, the question remains
open rather than being resolved
through an artificial harmonization
of the texts.


% =============================================================================
% 13. POINTS OF CONVERGENCE AND DIVERGENCE
% =============================================================================

\subsection{Points of Convergence and Divergence}
\label{subsec:ch08-farabi-convergence-divergence}

% Structural ID: VS.B01.P01.C08.S12.SS13.P001
\phantomsection\label{para:VS.B01.P01.C08.S12.SS13.P001}

A first point of convergence
concerns the refusal to treat
Being as an empty label
whose meaning is exhausted
by a single act of predication.
Al-Farabi's distinctions
among the senses of Being
show that metaphysical inquiry
must identify the respect
in which something is said
to be. Book I likewise
distinguishes Being from
Identity and does not
reduce either to a merely
descriptive or linguistic
function. The convergence
is methodological and
conceptual, not an identity
of doctrines.

% Structural ID: VS.B01.P01.C08.S12.SS13.P002
\phantomsection\label{para:VS.B01.P01.C08.S12.SS13.P002}

A second point of convergence
concerns the analysis of Unity.
Al-Farabi's treatment of
multiple senses of one
supports the requirement
that Unity be examined
according to its subject
and respect. Book I similarly
rejects the reduction of
Unity to Singularity,
Aggregate, or Composition.
The two projects therefore
share a concern with
distinguishing forms of
unity that might otherwise
be conflated.

% Structural ID: VS.B01.P01.C08.S12.SS13.P003
\phantomsection\label{para:VS.B01.P01.C08.S12.SS13.P003}

A third point of convergence
concerns the ontological
significance of specification.
Al-Farabi's account of
circumscription by quiddity
shows that the distinction
of a thing from what it is
not cannot be adequately
reduced to naming or
external description. Book I
also rejects the idea that
Identity is constituted
by recognition or linguistic
designation. The convergence
does not establish that
quiddity and Identity are
the same concept.

% Structural ID: VS.B01.P01.C08.S12.SS13.P004
\phantomsection\label{para:VS.B01.P01.C08.S12.SS13.P004}

A principal divergence concerns
the explanatory role of
the First Existent. Al-Farabi's
system gives the First a
unique causal and metaphysical
priority from which the
dependent order of existents
is explained. Book I does
not identify Identity with
a first cause, nor does
it derive the plurality
of identities from a single
causal principle. The
metaphysical priority of
Al-Farabi's First cannot
be transferred to the
foundational relation between
Being and Identity.

% Structural ID: VS.B01.P01.C08.S12.SS13.P005
\phantomsection\label{para:VS.B01.P01.C08.S12.SS13.P005}

A further divergence concerns
quiddity and Determination.
Al-Farabi's account of
circumscription by quiddity
is relevant to the specification
of an existent, but Book I
does not define Identity
as the sum of its determinations
or as a synonym for whatness.
The constitutive relation
between Identity and
Determination requires
a separate argument. The
presence of a determinate
quiddity cannot be treated
as a complete substitute
for that argument.

% Structural ID: VS.B01.P01.C08.S12.SS13.P006
\phantomsection\label{para:VS.B01.P01.C08.S12.SS13.P006}

The same restriction applies
to emanation and Manifestation.
Emanation is a causal account
of how dependent existents
proceed from the First.
Manifestation concerns
the ontological problem
examined in Book I and
cannot be reduced to the
causal production of an
existent by a first principle.
Likewise, the Active Intellect
and the celestial hierarchy
cannot be identified with
Light or with the proposed
domains and hubs of Book I.


% =============================================================================
% 14. THE STRONGEST CHALLENGE AND THE RESPONSE OF BOOK I
% =============================================================================

\subsection{The Strongest Challenge and the Response of Book I}
\label{subsec:ch08-farabi-strongest-challenge}

% Structural ID: VS.B01.P01.C08.S12.SS14.P001
\phantomsection\label{para:VS.B01.P01.C08.S12.SS14.P001}

The strongest challenge that
Al-Farabi's account presents
to Book I concerns the
relation between Being,
Unity, quiddity, and the
specification of an existent.
If a thing is one through
being circumscribed by its
quiddity, and if this sense
of one is coextensive with
a corresponding sense
of Being, then an ontology
of Identity must explain
why Identity remains
analytically distinguishable
from Determination and
Unity. It is not sufficient
to assert that these concepts
are different; their
ontological relation must
be clarified.
\citep[sec.~3.6]{MennAlFarabiMetaphysicsSEP2026}

% Structural ID: VS.B01.P01.C08.S12.SS14.P002
\phantomsection\label{para:VS.B01.P01.C08.S12.SS14.P002}

The challenge becomes stronger
when individual specification
is considered. If an existent
is distinguished from other
existents through a quiddity
that circumscribes it, then
Book I must explain whether
the Identity of that existent
is exhausted by this
specification or whether
quiddity is one dimension
of a more comprehensive
ontological constitution.
The question cannot be
dismissed by treating
quiddity as merely a
linguistic definition.

% Structural ID: VS.B01.P01.C08.S12.SS14.P003
\phantomsection\label{para:VS.B01.P01.C08.S12.SS14.P003}

The response of Book I begins
by preserving the analytical
distinction among Being,
Identity, Determination,
and Unity. Being concerns
the actuality of what exists.
Identity concerns the
existent as the existent
that it is. Determination
concerns the determinate
character of an existent
or of an aspect of it.
Unity concerns a specified
mode of being one. These
concepts may be ontologically
related, but their relations
must not be converted into
conceptual equivalences.

% Structural ID: VS.B01.P01.C08.S12.SS14.P004
\phantomsection\label{para:VS.B01.P01.C08.S12.SS14.P004}

The claim that a particular
sense of Unity is coextensive
with a particular sense
of Being does not establish
that Unity is identical
with Identity. Nor does
the fact that quiddity
specifies an existent
establish that Identity
is reducible to quiddity.
The argument concerns
the scope and function
of the concepts. A relation
of coextension within
a specified conceptual
analysis does not by itself
establish that the related
terms are identical in
every ontological respect.

% Structural ID: VS.B01.P01.C08.S12.SS14.P005
\phantomsection\label{para:VS.B01.P01.C08.S12.SS14.P005}

The response also distinguishes
causal dependence from
ontological Identity. That
an existent depends upon
the First Existent for
its existence does not
by itself explain the
specific Identity of the
dependent existent. The
causal explanation of
why an existent exists
and the ontological
question of what makes
it the existent that it
is are connected, but
they are not the same
question. An adequate
account must address
both without reducing
one to the other.

% Structural ID: VS.B01.P01.C08.S12.SS14.P006
\phantomsection\label{para:VS.B01.P01.C08.S12.SS14.P006}

The response does not establish
that Al-Farabi's analysis
is mistaken. It establishes
a limit on what can be
inferred from it. His
account of Being, Unity,
quiddity, and causal
dependence is a substantive
metaphysical contribution.
It does not, without
further argument, settle
the foundational relation
between Being and Identity
adopted in Book I. The
challenge therefore remains
a productive demand for
conceptual clarification,
not a reason to abandon
the distinctions of Part I.


% =============================================================================
% 15. BILATERAL COSTS AND THE KIND OF RELATION
% =============================================================================

\subsection{Bilateral Costs and the Kind of Relation}
\label{subsec:ch08-farabi-bilateral-costs}

% Structural ID: VS.B01.P01.C08.S12.SS15.P001
\phantomsection\label{para:VS.B01.P01.C08.S12.SS15.P001}

The comparison has a cost for
Book I. It cannot treat
Al-Farabi's analysis of
Being, Unity, and quiddity
as a historical curiosity
or as a preliminary version
of concepts already developed
in the present framework.
His distinctions reveal
that a theory of Identity
must explain the ontological
significance of specification,
the different senses of
Unity, and the relation
between a dependent existent
and its causal ground.

% Structural ID: VS.B01.P01.C08.S12.SS15.P002
\phantomsection\label{para:VS.B01.P01.C08.S12.SS15.P002}

The corresponding cost for
Al-Farabi's account, when
examined from the questions
of Book I, is that causal
dependence and quiddity
do not automatically settle
the Identity of every
dependent existent. His
system provides an account
of the First and the order
that proceeds from it,
but the question of how
Identity is to be distinguished
from quiddity, determination,
and Unity remains a
separate comparative
problem. This is a limit
relative to the inquiry
of Book I, not a claim
that Al-Farabi's own
system is thereby internally
incoherent.

% Structural ID: VS.B01.P01.C08.S12.SS15.P003
\phantomsection\label{para:VS.B01.P01.C08.S12.SS15.P003}

The relation between the
two projects is therefore
one of critical proximity
without conceptual equivalence.
They share concerns with
Being, Unity, specification,
and the ordered plurality
of existents. They differ
in their explanatory
architecture and in the
role assigned to the
First Existent, quiddity,
causal dependence, and
Identity. Neither system
can be substituted for
the other without changing
the questions being asked.

% Structural ID: VS.B01.P01.C08.S12.SS15.P004
\phantomsection\label{para:VS.B01.P01.C08.S12.SS15.P004}

The substitution test confirms
this limit. Replacing
Identity in Book I with
Al-Farabi's First Existent
would change the meaning
of the claims: the First
Existent is a first cause
within a metaphysical
order, whereas Identity
concerns the existent
as the existent that
it is. Replacing Identity
with quiddity would
likewise change the
argument, because quiddity
specifies what a thing
is, while Identity is
not defined as a synonym
for whatness.

% Structural ID: VS.B01.P01.C08.S12.SS15.P005
\phantomsection\label{para:VS.B01.P01.C08.S12.SS15.P005}

Replacing Manifestation
with emanation would
also change the explanatory
function. Emanation concerns
the causal derivation
of dependent beings from
the First. Manifestation
belongs to a distinct
ontological inquiry in
Book I. Replacing Light
with the Active Intellect
would similarly introduce
a separate intellectual
substance and a specific
cosmological and
psychological function
that cannot be assumed
within the present
framework.

% Structural ID: VS.B01.P01.C08.S12.SS15.P006
\phantomsection\label{para:VS.B01.P01.C08.S12.SS15.P006}

The relation is consequently
not identity, equivalence,
or direct doctrinal
continuity. It is a
critical comparison
between two distinct
metaphysical projects
that address overlapping
problems while preserving
different conceptual
commitments. The value
of the comparison lies
in the clarification
of these overlaps and
differences, not in
the construction of
a synthetic system
that erases them.


% =============================================================================
% 16. OPEN PROBLEMS AND STATUS OF CLAIMS
% =============================================================================

\subsection{Open Problems and Status of Claims}
\label{subsec:ch08-farabi-open-problems}

% Structural ID: VS.B01.P01.C08.S12.SS16.P001
\phantomsection\label{para:VS.B01.P01.C08.S12.SS16.P001}

The comparison leaves several
questions open. First, what
is the exact relation between
Al-Farabi's account of
quiddity and the ontological
problem of Identity? Second,
does the coextension of
Being and Unity in a
specified sense have
implications for the
foundational relation
between Being and Identity
in Book I? Third, how
should causal dependence
be related to the
determinate constitution
of dependent existents?
These questions require
further argument and
cannot be settled through
terminological resemblance.

% Structural ID: VS.B01.P01.C08.S12.SS16.P002
\phantomsection\label{para:VS.B01.P01.C08.S12.SS16.P002}

A further open problem
concerns the relation
between a first source
and the plurality of
existents. Al-Farabi's
account explains the
asymmetrical dependence
of the many upon the
First. Book I must
clarify whether and
how the notion of a
source relates to
Identity without
reducing the plurality
of identities to a
single identity or
to an undifferentiated
aggregate of effects.

% Structural ID: VS.B01.P01.C08.S12.SS16.P003
\phantomsection\label{para:VS.B01.P01.C08.S12.SS16.P003}

The table below classifies
the principal claims of
this section along two
independent axes. The
first identifies the
functional or logical
form of each claim.
The second identifies
its epistemic or
justificatory status
within the present
comparative inquiry.
These classifications
describe the claims
as reconstructed here;
they are not presented
as Al-Farabi's own
classification of
knowledge or argument.

% =============================================================================
% Project : THE PHILOSOPHY OF VIABILITY
% Book    : Book I — Ontology of Viability
% Book ID : VS.B01
% Part    : Part I — Identity and Manifestation
% Part ID : VS.B01.P01
% Chapter : Chapter 08 — Comparative Philosophical Examination
% Chapter ID : VS.B01.P01.C08
%
% File :
%   4_ASSETS/tables/
%   VS.B01.P01.C08.S12.SS16.P003.TL01.al_farabi_claims_and_comparative_limits.tex
%
% Section : Section 12 — Al-Farabi
%
% Asset :
%   Table TL01 — Al-Farabi's Claims and Comparative Limits
%
% Structural Identifier :
%   VS.B01.P01.C08.S12.SS16.P003.TL01
%
% Purpose :
%   To distinguish the historical reconstruction of Al-Farabi's
%   principal metaphysical claims from their logical form,
%   epistemic status within the present inquiry, and comparative
%   limits for the Ontology of Viability.
%
% Intellectual Benefit :
%   Prevents the conflation of historical doctrine, interpretive
%   reconstruction, and the foundational commitments of Book I.
%
% Aesthetic / Experiential Benefit :
%   Provides a compact reference for the extended comparative
%   analysis in Section 12.
%
% Representational Limit :
%   The classifications concern the present reconstruction.
%   They are not Al-Farabi's own taxonomy of claims or knowledge.
%
% Dependencies :
%   longtable, booktabs, array.
%
% =============================================================================

\small
\renewcommand{\arraystretch}{1.20}
\setlength{\tabcolsep}{4pt}

\begin{longtable}{
    >{\raggedright\arraybackslash}p{2.0cm}
    >{\raggedright\arraybackslash}p{2.6cm}
    >{\raggedright\arraybackslash}p{2.25cm}
    >{\raggedright\arraybackslash}p{2.55cm}
    >{\raggedright\arraybackslash}p{4.0cm}
}

\caption{Al-Farabi's Principal Claims and Their Comparative Limits}
\label{tab:VS.B01.P01.C08.S12.SS16.P003.TL01} \\

\toprule
\textbf{Claim or issue}
&
\textbf{Historical reconstruction}
&
\textbf{Functional or logical form}
&
\textbf{Epistemic or justificatory status}
&
\textbf{Comparative limit for Book I}
\\
\midrule
\endfirsthead

\toprule
\textbf{Claim or issue}
&
\textbf{Historical reconstruction}
&
\textbf{Functional or logical form}
&
\textbf{Epistemic or justificatory status}
&
\textbf{Comparative limit for Book I}
\\
\midrule
\endhead

\midrule
\multicolumn{5}{r}{\footnotesize\textit{Continued on next page}}\\
\endfoot

\bottomrule
\endlastfoot

Being
&
Al-Farabi distinguishes several senses of Being, including Being as true and Being as circumscribed by a quiddity outside the soul.
&
Proposition
&
Argued or Derived Claims
&
The differentiation of senses does not establish that Being is identical with Identity or that its ontological status depends upon cognition.
\\

\addlinespace

Unity
&
Al-Farabi distinguishes numerical unity, continuity, indivisibility, uniqueness, and circumscription by quiddity.
&
Proposition
&
Argued or Derived Claims
&
The plurality of senses of Unity supports conceptual differentiation but does not establish that Unity is identical with Identity, Singularity, or Composition.
\\

\addlinespace

Unity and quiddity
&
A particular sense of being one is connected with the circumscription of an existent by its quiddity.
&
Proposition
&
Conditionally Supported Claims
&
A specified relation between Unity and quiddity cannot be generalized into an identity of Unity, quiddity, and Identity.
\\

\addlinespace

The First Existent
&
The First is presented as uncaused, independent, simple, unique, and the first cause of the existence of other beings.
&
Foundational Commitment within Al-Farabi's system
&
Conditionally Supported Claims
&
The First Existent is not equivalent to Being as such, Identity, or any proposed realm of Book I.
\\

\addlinespace

Emanation
&
Dependent existents proceed from the First without adding to its perfection or requiring an alteration in it.
&
Thesis
&
Conditionally Supported Claims
&
Causal production is not equivalent to Expression or Manifestation, and causal dependence does not by itself explain the Identity of the dependent existent.
\\

\addlinespace

Secondary intellects and celestial order
&
The First is followed by an ordered series of separate intellects associated with the celestial structure; details differ between relevant treatises.
&
Proposition
&
Conditionally Supported Claims
&
The cosmological hierarchy is not equivalent to the proposed domains and hubs of Book I.
\\

\addlinespace

Active Intellect
&
The Active Intellect occupies a distinct position in the separate-intellect order and is connected with the sublunary domain and human intellectual activity.
&
Proposition
&
Argued or Derived Claims
&
The Active Intellect is not identical with Light or with a foundational component of the ontology of Book I.
\\

\addlinespace

The virtuous city
&
The city is related analogically to the ordered cosmos and to the differentiated organization of an organism.
&
Thesis
&
Conditionally Supported Claims
&
The analogy does not by itself establish the normative legitimacy of a political hierarchy or define the Identity of a social entity.
\\

\addlinespace

Relation to Book I
&
Al-Farabi's analysis of Being, Unity, quiddity, and causal dependence overlaps with problems examined in Part I but uses a distinct explanatory architecture.
&
Thesis
&
Argued or Derived Claims
&
The relation is critical proximity without conceptual equivalence or direct doctrinal continuity.
\\

\addlinespace

Identity and quiddity
&
The relation between Al-Farabi's account of circumscription by quiddity and the ontological status of Identity remains open to further analysis.
&
Open Problem
&
Unresolved or Open Problems
&
The question cannot be settled by substituting Identity for quiddity or by treating the two terms as interchangeable.
\\

\end{longtable}

\normalsize

% =============================================================================
% END OF TABLE TL01
% =============================================================================

% Structural ID: VS.B01.P01.C08.S12.SS16.P004
\phantomsection\label{para:VS.B01.P01.C08.S12.SS16.P004}

The table does not place
the functional form of
a claim and its epistemic
status on a single scale.
A proposition may be
conditionally supported;
a thesis may be argued;
and an open problem may
be precisely formulated
without being resolved.
The two axes therefore
remain independent.
The classification also
prevents a historical
reconstruction from
being mistaken for
a foundational commitment
of Book I.

% Structural ID: VS.B01.P01.C08.S12.SS16.P005
\phantomsection\label{para:VS.B01.P01.C08.S12.SS16.P005}

The present reading would
require revision if a
more complete examination
of the primary texts
established that Al-Farabi
explicitly identifies
quiddity, Unity, and
Identity in the same
technical sense adopted
in Book I, or if the
differences between
the relevant treatises
were shown to have
been misunderstood.
Conversely, the comparative
thesis would be weakened
if it could not explain
why the concepts of
Being, quiddity, Unity,
and Identity perform
different explanatory
roles in the texts
under examination.


% =============================================================================
% 17. TRANSITION TO AVICENNA
% =============================================================================

\subsection{Transition to Avicenna}
\label{subsec:ch08-farabi-transition}

% Structural ID: VS.B01.P01.C08.S12.SS17.P001
\phantomsection\label{para:VS.B01.P01.C08.S12.SS17.P001}

Al-Farabi's metaphysics has
shown that the analysis
of Being and Unity cannot
be separated from the
questions of quiddity,
causal dependence, and
the ordered plurality
of existents. His account
of the First Existent
provides a distinctive
explanation of the
causal order of reality,
while his analysis of
Unity reveals several
ways in which a thing
may be one and
distinguished from
what it is not.

% Structural ID: VS.B01.P01.C08.S12.SS17.P002
\phantomsection\label{para:VS.B01.P01.C08.S12.SS17.P002}

The comparison has also
shown that these concepts
must not be treated as
synonyms for Identity.
The First Existent,
quiddity, Unity, and
emanation perform
different philosophical
functions. Their relation
to Being and Identity
remains a substantive
ontological problem,
not a matter of
translating one
vocabulary into another.

% Structural ID: VS.B01.P01.C08.S12.SS17.P003
\phantomsection\label{para:VS.B01.P01.C08.S12.SS17.P003}

The next section turns
to Ibn Sina. It examines
his more systematic
account of essence
and existence, the
Necessary Existent,
causal dependence,
and the hierarchy
of beings. The transition
must preserve the
historical distinction
between Al-Farabi's
metaphysical analyses
and Avicenna's later
formulation. The
questions are related,
but their answers
must be reconstructed
from their respective
texts rather than
assumed to be identical.


% =============================================================================
% END OF SECTION 12 OF CHAPTER 08
% =============================================================================